\documentclass[11pt]{article}
\usepackage[margin=2.5cm]{geometry}
\usepackage{amsmath,amssymb,bm}
\usepackage{booktabs,array}
\usepackage{hyperref}
\usepackage{enumitem}
\setlist{itemsep=2pt,topsep=4pt}

\newcommand{\ext}{\mathrm{ext}}
\newcommand{\pd}[2]{\frac{\partial #1}{\partial #2}}

\title{Open boundary conditions for the barotropic transport:\\
Flather, Dirichlet, and a balanced/unbalanced split}
\author{Notes for the PanArcticShenanigans boundary experiments}
\date{31 October 2026}

\begin{document}
\maketitle

\begin{abstract}
These notes set out the problem in the linear hydrostatic equations, reduce it to rotating shallow water through
vertical modes, and explain the open boundary conditions to be compared: Flather, Dirichlet, Davies relaxation,
Flather with a pinned net transport, and a ``balanced Flather'' condition in which the external data prescribe the
balanced (geostrophic) part of the flow and only the unbalanced part is radiated. The balanced part is defined
from potential vorticity (PV) inversion, with the streamfunction on the open boundary set by the external
transport. The last section describes a test: a basin with three walls and one open side, crossed by a steady
Gaussian vortex.
\end{abstract}

%=====================================================================================================
\section{Linear hydrostatic equations and vertical modes}

Linearise the hydrostatic Boussinesq equations on an $f$-plane about a state of rest with stratification
$N^2(z)$, flat bottom at $z=-H$ and a free surface $\eta$. With $p$ the kinematic pressure perturbation and $b$ the
buoyancy perturbation,
\begin{align}
  u_t - f v &= -p_x, & v_t + f u &= -p_y, & p_z &= b, \\
  b_t + N^2 w &= 0,  & u_x + v_y + w_z &= 0, &&
\end{align}
with $w=0$ at $z=-H$ and, linearised at $z=0$, $w=\eta_t$ and $p=g\eta$.

\paragraph{Separation.} Write $(u,v,p) = \sum_n (u_n,v_n,p_n)(x,y,t)\,F_n(z)$. The vertical structure solves
\begin{equation}
  \frac{d}{dz}\!\left(\frac{1}{N^2}\frac{dF_n}{dz}\right) + \frac{F_n}{c_n^2} = 0,\qquad
  F_n'(-H) = 0,\qquad F_n'(0) + \frac{N^2(0)}{g}F_n(0) = 0,
  \label{eq:sturm}
\end{equation}
and each mode obeys the rotating shallow-water equations with its own speed $c_n$:
\begin{equation}
  \partial_t u_n - f v_n = -\partial_x p_n,\qquad
  \partial_t v_n + f u_n = -\partial_y p_n,\qquad
  \partial_t p_n + c_n^2\,(\partial_x u_n + \partial_y v_n) = 0.
  \label{eq:swmode}
\end{equation}
The barotropic mode has $c_0\approx\sqrt{gH}$ and $F_0\approx 1$ (exactly so with a rigid lid); baroclinic modes
have $c_n = O(1\text{ m s}^{-1})$. Everything below is written for one mode; the barotropic case is the one that
sets the transport.

\paragraph{Barotropic mode in transport form.} With $U=Hu$, $V=Hv$ and $c=\sqrt{gH}$:
\begin{equation}
  U_t - fV = -gH\,\eta_x,\qquad V_t + fU = -gH\,\eta_y,\qquad \eta_t + U_x + V_y = 0.
  \label{eq:sw}
\end{equation}
This is what a hydrostatic model with $N_z=1$ and no advection solves, and what the split-explicit free surface
solves for the depth-integrated transport in the full model.

%=====================================================================================================
\section{Waves, balance and potential vorticity}

\paragraph{Normal modes.} Plane waves $e^{i(kx+ly-\omega t)}$ of \eqref{eq:sw} satisfy
\begin{equation}
  \omega\left(\omega^2 - f^2 - c^2(k^2+l^2)\right) = 0 .
\end{equation}
There are two kinds of motion:
\begin{itemize}
  \item the \emph{geostrophic} (balanced) mode, $\omega=0$: steady, with $fv = g\eta_x$, $fu=-g\eta_y$;
  \item \emph{inertia--gravity} (unbalanced) waves, $|\omega|\ge f$: dispersive, with phase and group speeds that
        depend on wavenumber; non-dispersive at speed $c$ only for $\omega\gg f$.
\end{itemize}
With $\beta$ or topography the balanced mode becomes slow Rossby or topographic waves rather than strictly steady,
but the frequency gap at $f$ remains.

\paragraph{Potential vorticity.} The linear PV
\begin{equation}
  q = \zeta - \frac{f}{H}\,\eta,\qquad \zeta = v_x - u_y,
\end{equation}
is conserved pointwise: $q_t = 0$. Inertia--gravity waves carry no PV; all PV is in the balanced part.

\paragraph{Balanced/unbalanced decomposition.} The balanced part is the geostrophic state with the same PV. With
$\psi = g\eta/f$ (so $u=-\psi_y$, $v=\psi_x$ for geostrophic flow),
\begin{equation}
  \left(\nabla^2 - \frac{1}{L_d^2}\right)\psi_B = q,\qquad L_d = \frac{c}{f},
  \label{eq:inversion}
\end{equation}
and
\begin{equation}
  \eta_B = \frac{f}{g}\psi_B,\quad u_B=-\partial_y\psi_B,\quad v_B=\partial_x\psi_B;\qquad
  (\eta,u,v) = (\eta_B,u_B,v_B) + (\eta',u',v').
\end{equation}
The remainder $(\eta',u',v')$ has $q'=0$ and is a sum of inertia--gravity waves. On an unbounded or periodic
$f$-plane the two parts are orthogonal in the energy norm $E=\tfrac12\int\left(H(u^2+v^2)+g\eta^2\right)dA$, so the
energy splits cleanly into balanced and unbalanced parts.

For a baroclinic mode, \eqref{eq:inversion} holds with $c_n$, so $L_{d,n}$ is a few kilometres in the Arctic: the
inversion is then nearly local. For the barotropic mode $L_d$ is $O(1000\text{ km})$ and the inversion is global.

%=====================================================================================================
\section{Boundary conditions for the streamfunction}
\label{sec:psibc}

Equation \eqref{eq:inversion} needs $\psi_B$ on the whole boundary.

\begin{itemize}
  \item \textbf{Walls.} No normal flow: $\psi_B$ is constant along a connected coastline, say $\psi_B=\psi_w$.
  \item \textbf{Open boundary} (south side, $y=0$, $x_W\le x\le x_E$). The normal velocity of the balanced flow is
        $v_B = \partial_x\psi_B$. If the balanced normal flow is the external one, then
        \begin{equation}
          \psi_B(x,0) = \psi_w + \int_{x_W}^{x} v_\ext(x')\,dx'.
          \label{eq:psiopen}
        \end{equation}
        Continuity at the south-east corner requires $\int_{x_W}^{x_E} v_\ext\,dx = 0$: a steady balanced state in a
        basin with one connected coastline has zero net transport through the open side. A non-zero net transport
        changes the basin volume and is treated separately (Section~\ref{sec:net}).
\end{itemize}

So the external transport \emph{is} the Dirichlet condition for the streamfunction on the open side, and the
balanced state is fixed up to the single constant $\psi_w$.

\paragraph{Fixing the constant.} By linearity, $\psi_B = \psi_p + \psi_w\,\phi_1$, where $\psi_p$ solves
\eqref{eq:inversion} with $\psi_w=0$ in \eqref{eq:psiopen} and $\phi_1$ solves the homogeneous problem with
$\phi_1=1$ on the whole boundary (computed once). Requiring the balanced part to carry the basin volume,
$\int \eta_B\,dA = \int \eta\,dA$, gives
\begin{equation}
  \psi_w = \frac{\tfrac{g}{f}\int\eta\,dA - \int\psi_p\,dA}{\int\phi_1\,dA}.
  \label{eq:psiw}
\end{equation}

\paragraph{The balanced free surface on the open side.} From \eqref{eq:psiopen}, $\eta_B$ along the open side has
the same shape as the external free surface, which is itself geostrophic with $v_\ext$:
\begin{equation}
  \eta_B(x,0,t) = \eta_\ext(x,t) + C(t),
  \label{eq:etaB}
\end{equation}
where $C$ is a single constant set by \eqref{eq:psiw}. This is the key relation: \emph{the model's balanced free
surface may differ from the external one along the open side only by a constant.}

%=====================================================================================================
\section{One-dimensional characteristics and reflection}

Take \eqref{eq:sw} without rotation and without $x$-dependence near the south boundary, domain $y>0$:
$\eta_t+V_y=0$, $V_t+c^2\eta_y=0$. The Riemann invariants are
\begin{equation}
  R^+ = V + c\eta \;\;(\text{speed } +c,\ \text{into the domain}),\qquad
  R^- = V - c\eta \;\;(\text{speed } -c,\ \text{out of the domain}).
\end{equation}
At the boundary $R^-$ is known from the interior; one condition is needed, on $R^+$. Consider the family
\begin{equation}
  V + \alpha\,c\,\eta = \text{data},\qquad \alpha\ge 0 .
\end{equation}
Substituting $V=(R^++R^-)/2$, $c\eta=(R^+-R^-)/2$, an outgoing wave $R^-$ is reflected into
\begin{equation}
  R^+_{\text{refl}} = r\,R^-,\qquad r = -\frac{1-\alpha}{1+\alpha}.
\end{equation}
Flather is $\alpha=1$ ($r=0$); Dirichlet on the transport is $\alpha=0$ ($r=-1$); Dirichlet on $\eta$ is
$\alpha\to\infty$ ($r=+1$).

\paragraph{With rotation.} Coriolis terms are lower order, so the characteristic structure (one incoming
characteristic at speed $c$, one outgoing, one at zero speed for the tangential flow) is unchanged, and so is the
number of conditions. What rotation changes is the quality of absorption: Flather is nearly transparent for
$\omega\gg f$ at normal incidence and only partly so near $\omega=f$ or at oblique incidence. This is measurable in
the test below.

\paragraph{Why Flather turns a free-surface offset into transport.} For a balanced flow, $V$ and $\eta$ along the
boundary are linked by geostrophy, $V=(gH/f)\,\eta_x$, not by $V=\pm c\eta$. A constant offset of the model's
balanced $\eta$ from $\eta_\ext$ is not a wave, but Flather reads it as one and adds a transport $c\,C$ per unit
length. With $c=170\text{ m s}^{-1}$, $C=3\text{ cm}$ and a 4000~km boundary this is about $20$~Sv.

%=====================================================================================================
\section{The boundary conditions compared}

All conditions act on the normal barotropic transport on the open faces, at every barotropic substep. On the
C-grid, $V$ lives on the boundary face, $\eta$ at cell centres; $\eta_b$ is the face average of the first
interior and the halo value. Written for the south side; for a north side the sign of the $c$ term flips.

\begin{description}
  \item[M1 Flather.] $V_b = V_\ext - c\,(\eta_b - \eta_\ext)$.
        Imposes the incoming characteristic $V+c\eta$. Transport depends on the $\eta$ mismatch.

  \item[M2 Dirichlet.] $V_b = V_\ext$.
        Transport exact; every outgoing barotropic wave reflects ($r=-1$). In a linear inviscid model reflected
        energy stays in the basin.

  \item[M3 Davies relaxation.] M2 on the face, plus $\partial_t\phi \mathrel{+}= -\gamma(y)\,(\phi-\phi_\ext)$ for
        $\phi\in\{U,V,\eta\}$ in a rim of width $W$, with $\gamma$ decreasing smoothly from $\gamma_{\max}$ at the
        boundary to $0$ at $y=W$. Absorbs waves whose wavelength is not much larger than $W$. For the barotropic
        mode $\lambda\sim 2\pi c/\omega$ is thousands of kilometres, so the rim would have to be very wide; for
        baroclinic modes $\lambda$ is tens of kilometres and a rim works. The rim also relaxes the balanced part to
        the external state.

  \item[M4 Flather with pinned net transport] (current \texttt{target\_transport}).
        $V_b = V_\ext - c(\eta_b-\eta_\ext) - S$, with $S$ uniform along the boundary and chosen so that
        $\int V_b\,dx = \int V_\ext\,dx$. The correction is uniform \emph{transport} per unit length, so the
        velocity correction scales as $1/H$: largest on shelves.

  \item[M5 Balanced Flather.] External data prescribe the balanced part; only the unbalanced part radiates, and
        no unbalanced wave enters from outside:
        \begin{equation}
          V'_b + c\,\eta'_b = 0
          \quad\Longrightarrow\quad
          V_b = V_{B,b} - c\,(\eta_b - \eta_{B,b}) = V_\ext - c\,\bigl(\eta_b - \eta_\ext - C\bigr),
          \label{eq:m5}
        \end{equation}
        using $V_{B,b}=V_\ext$ and \eqref{eq:etaB}. It differs from Flather only by the constant $C$.
        \begin{description}
          \item[M5a] $C$ from the PV inversion, \eqref{eq:inversion}--\eqref{eq:psiw}. Unbalanced waves that change
                the basin volume can leave.
          \item[M5b] $C$ from the net constraint $\int V_b\,dx = \int V_\ext\,dx$:
                \begin{equation}
                  C = \frac{\int c\,(\eta_b-\eta_\ext)\,dx}{\int c\,dx}.
                \end{equation}
                No inversion; equivalent to M4 with the correction weighted by $c=\sqrt{gH}$ instead of uniform,
                so the velocity correction scales as $1/\sqrt{H}$. The net unbalanced flux is forced to zero, so the
                volume-changing part of an outgoing wave reflects.
        \end{description}
\end{description}

\begin{table}[h]
\centering
\begin{tabular}{@{}l>{\raggedright}p{3.6cm}>{\raggedright}p{3.4cm}>{\raggedright\arraybackslash}p{4.4cm}@{}}
\toprule
 & Transport set by & Outgoing waves & Expected weakness \\
\midrule
M1 Flather          & $\eta$ mismatch          & leave ($\omega\gg f$)   & offset $\Rightarrow$ $c\,C$ per length \\
M2 Dirichlet        & external data            & reflect                 & trapped wave energy \\
M3 Davies           & external data            & damped in rim           & rim width for barotropic waves; cost \\
M4 Flather + pin    & external net             & leave, except net part  & $1/H$ correction drains shelves \\
M5a balanced (PV)   & external (balanced)      & leave                   & inversion cost; $\beta$, topography \\
M5b balanced (net)  & external net             & leave, except net part  & volume mode reflects \\
\bottomrule
\end{tabular}
\caption{Summary of the conditions on the barotropic transport.}
\end{table}

%=====================================================================================================
\section{Net transport and basin volume}
\label{sec:net}

If the external net transport $Q(t)=\int V_\ext\,dx$ is non-zero, the basin volume changes and $\int\eta\,dA$
follows $Q$. In the decomposition this is a spatially uniform change in $\eta$, carried by the constant $\psi_w$
(equivalently $C$); it has no geostrophic flow. M2, M3, M4 and M5b impose $Q$ exactly; M1 imposes it only if the
mean edge mismatch vanishes; M5a imposes the balanced part and lets the unbalanced part carry its own net flux.
For a long run, an error $\delta Q$ in the mapped external net transport changes the mean free surface by
$\delta Q/A$ per unit time (about $0.5$ cm/day per Sv for the pan-Arctic domain), whichever condition is used.

%=====================================================================================================
\section{From modes to depth in the stratified model}

For each open-boundary column, compute the modes $F_n(z)$ from that column's $N^2(z)$ by \eqref{eq:sturm}, once.
Each step:
\begin{enumerate}
  \item project the model and external normal-velocity profiles onto the modes,
        $v(z)=\sum_n a_n F_n(z)$ (an $N_z\times N_z$ matrix--vector product per column);
  \item apply the chosen condition to each $a_n$ with its own $c_n$ (and, for M5, its own balanced part);
  \item rebuild $v_b(z)=\sum_n a_n^b F_n(z)$.
\end{enumerate}
The cost is negligible: a few hundred boundary columns. The split-explicit model already treats the depth integral
(close to mode 0) separately, so the first step is to change only the barotropic condition; mode-by-mode
baroclinic conditions are a later refinement.

%=====================================================================================================
\section{Test problem}

\paragraph{Model.} Oceananigans \texttt{HydrostaticFreeSurfaceModel}, $N_z=1$, \texttt{SplitExplicitFreeSurface},
no momentum advection, no buoyancy, $f$-plane. This is \eqref{eq:sw} with the same open-boundary machinery as the
pan-Arctic model.

\paragraph{Geometry.} A square basin $0\le x\le L$, $0\le y\le L$, walls on the west, north and east, open on the
south. Suggested values: $L=2000$~km, $\Delta x=20$~km, $H=500$~m ($c=70\text{ m s}^{-1}$),
$f=1.4\times10^{-4}\text{ s}^{-1}$, so $L_d=500$~km. A shallower $H$ than the ocean keeps $L_d$ comparable to the
basin and slows the waves, so the reference run below can be affordable; the results depend on $R/L_d$ and
$L/L_d$.

\paragraph{Steady balanced state.} A Gaussian vortex, as in BalancedFlows.jl,
\begin{equation}
  \eta_v = \eta_0\,\exp\!\left(-\frac{(x-x_0)^2+(y-y_0)^2}{R^2}\right),\qquad
  u_v = -\frac{g}{f}\partial_y\eta_v,\quad v_v = \frac{g}{f}\partial_x\eta_v,
\end{equation}
centred on the open side ($x_0=L/2$, $y_0=0$), with $R=150$~km and $\eta_0=0.1$~m. The vortex flow crosses the open
side inwards on one half and outwards on the other, with zero net transport, and has decayed at the walls
($R\ll L/2$). It is an exact steady solution of \eqref{eq:sw}. The external data are $V_\ext=Hv_v$ and
$\eta_\ext=\eta_v$ on $y=0$.

\paragraph{Reference solution.} The same initial state in a domain extended far to the south (vortex fully inside,
the extra boundaries far enough that waves do not return within the run). Compare the two runs in the common
region.

\paragraph{Experiments.}
\begin{description}
  \item[E0 Steady.] Start from the vortex. All conditions should keep it steady; checks the implementation.
  \item[E1 Adjustment.] Vortex plus an unbalanced Gaussian bump in $\eta$ (no velocity) in the interior, e.g.
        $0.2$~m, $200$~km. Its PV sets a balanced remainder; the rest leaves as inertia--gravity waves.
        Tests reflection.
  \item[E2 Datum offset.] Model starts from the vortex; external data are $\eta_\ext=\eta_v+\delta$ with
        $\delta=3$~cm. The true solution is the steady vortex. This is the GLORYS problem in miniature.
  \item[E3 Slow net transport.] Add a slowly varying net inflow $Q(t)$ (period of tens of days) with the matching
        uniform rise of $\eta_\ext$. Tests whether slow transport is imposed.
  \item[E4 Shelf.] Repeat E1--E2 with a shelf along part of the open side ($H$ from 500~m to 50~m). Tests the
        distribution of the correction and coastal-trapped waves.
\end{description}

\paragraph{Diagnostics.}
\begin{itemize}
  \item Net transport through the open side and its distribution along $x$, against $V_\ext$.
  \item Balanced and unbalanced energy in the basin over time, from the decomposition of Section~2 applied to the
        model state: the unbalanced energy should leave; the balanced part should match the reference.
  \item Domain-mean $\eta$ and the constant $C(t)$.
  \item Maps of the error against the reference at chosen times.
\end{itemize}

\paragraph{Expected outcomes.}
E2 separates M1 (spurious transport $\approx c\,\delta\,\times$ width) from the rest; E1 separates M2 (trapped
waves) from M1, M5a and M5b; E4 separates M4 ($1/H$) from M5b ($1/\sqrt H$); M5a against M5b shows whether letting
the volume-changing wave part leave matters.

%=====================================================================================================
\section*{References}
\begin{itemize}
  \item Blayo, E. and Debreu, L. (2005). Revisiting open boundary conditions from the point of view of
        characteristic variables. \emph{Ocean Modelling} 9, 231--252.
  \item Davies, H. C. (1976). A lateral boundary formulation for multi-level prediction models.
        \emph{Q. J. R. Meteorol. Soc.} 102, 405--418.
  \item Flather, R. A. (1976). A tidal model of the north-west European continental shelf.
        \emph{Mém. Soc. R. Sci. Liège} 6(10), 141--164.
  \item Marchesiello, P., McWilliams, J. C. and Shchepetkin, A. (2001). Open boundary conditions for long-term
        integration of regional oceanic models. \emph{Ocean Modelling} 3, 1--20.
  \item Oddo, P. and Pinardi, N. (2008). Lateral open boundary conditions for nested limited area models: a scale
        selective approach. \emph{Ocean Modelling} 20, 134--156.
\end{itemize}

\end{document}
